\documentclass[11pt]{article}
\usepackage[letterpaper,hmargin=1.06in,vmargin=0.94in]{geometry}
\usepackage[T1]{fontenc}
\usepackage{lmodern,amsmath,amssymb,amsthm,microtype}
\usepackage[hidelinks]{hyperref}
\newtheorem{theorem}{Theorem}
\newtheorem{lemma}[theorem]{Lemma}
\newtheorem{corollary}[theorem]{Corollary}
\newcommand{\NL}{\mathsf{NL}}
\newcommand{\UL}{\mathsf{UL}}
\newcommand{\F}{\mathbb F}
\newcommand{\dist}{\operatorname{dist}}
\hypersetup{pdftitle={Masked binary words for unambiguous reachability}}
\title{Masked binary words for unambiguous reachability}
\author{}
\date{November 2026}
\begin{document}
\maketitle
\begin{abstract}
We give a proposed proof of $\NL=\UL$. The construction assigns each
vertex of a layered graph a binary word and applies the same random mask
to all the words. Each coordinate at which two candidate midpoint
vertices differ can force a bit of the mask if their minimum path costs
tie. A random binary Toeplitz matrix supplies the words. The construction
uses $O(\log N)$ random bits and polynomially bounded integer weights on
a graph with $N$ vertices.
\end{abstract}

\section{The weight assignment}

A simple layered directed graph has layers $0,\ldots,d$, and each edge
advances exactly one layer. An assignment of vertex weights is
\emph{min-isolating} if every reachable ordered pair has a unique
minimum-weight path. When comparing paths with fixed endpoints, we charge
only their internal vertices. Adding either endpoint to every path in a
comparison preserves all ties and minimizers.

\begin{theorem}\label{thm:main}
There is a logspace-computable family of nonnegative integer vertex
weights with the following guarantee. On each fixed simple layered
directed graph with $N\ge2$ vertices and $d\le N$, a uniform seed of
$O(\log N)$ bits gives a min-isolating assignment with probability at
least $1-N^{-2}$. Every path cost is bounded by $N^{O(1)}$.
\end{theorem}

Give the vertices distinct labels $x(v)\in\F_2^L$, where
$L=\lceil\log_2 N\rceil$. Set $k=16L$ and $B=16$. Use independent
uniform bits $\tau_0,\ldots,\tau_{k+L-2}$ to define
\[
    a_j(v)=\bigoplus_{i=0}^{L-1}\tau_{j-1+i}\,x_i(v)
                       \qquad(1\le j\le k).
\]
This is binary Toeplitz hashing with the order of the input coordinates
reversed. Choose $k$ further independent uniform bits
$\eta_1,\ldots,\eta_k$. For a vertex in positive layer $t$, set
\begin{equation}\label{eq:weight}
    w(v)=B^{\nu_2(t)}\sum_{j=1}^k B^{k-j}\eta_j a_j(v).
\end{equation}
Here $\nu_2(t)$ is the exponent of the highest power of two dividing
$t$. Set the weight in layer zero to zero. The two sources of bits
have different roles: $\tau$ assigns the words, then $\eta$ masks
the same positions in every word. The analysis of the mask applies
to \emph{arbitrary fixed} words.

\section{A bound for two competing midpoint vertices}

Fix endpoints $s,t$ with at least one internal layer, and consider their
open interval of layers. It has a unique layer $t_*$ of maximum binary
valuation $h$: between any two distinct integers of the same valuation
lies an integer of higher valuation. Call $t_*$ the pivot layer.
The entire interval lies strictly between consecutive multiples of
$2^{h+1}$, so the number of its layers of valuation $h-r$ is at most
$2^r$ for each $1\le r\le h$. In particular,
\begin{equation}\label{eq:geometry}
 \sum_{\substack{t'\text{ internal}\\t'\ne t_*}} B^{\nu_2(t')-h}
                 \ \le\ \sum_{r\ge1}(2/B)^r=\frac17 .
\end{equation}

For a pivot vertex $v$ lying on an $s$--$t$ path, let $D_v$ denote the
least full cost of an $s$--$t$ path going through $v$. The vertex $v$ is
prescribed, and all other internal vertices may vary in taking this minimum.

\begin{lemma}\label{lem:mask}
Fix the entire word table $a$. For any two distinct feasible vertices
$u,v$ on the pivot layer,
\[
           \Pr_\eta[D_u=D_v\mid a]\le
                    2^{-\dist(a(u),a(v))},
\]
where $\dist$ denotes Hamming distance.
\end{lemma}

\begin{proof}
Fix a coordinate $j$ and expose $\eta_1,\ldots,\eta_{j-1}$.
Let $c_u$ be the minimum cost over paths through $u$, charging the
contributions of coordinates $1,\ldots,j-1$ at all internal vertices.
Define $c_v$ the same way. Both numbers are determined by the exposed
prefix. Write $\sigma=B^{h+k-j}$ for one unit in coordinate $j$ at
the pivot. On every path through $u$, the current coordinate at the
pivot contributes exactly $\sigma\eta_j a_j(u)$.

Two contributions remain. At the current coordinate, the lower layers
together cost at most $\sigma/7$ by \eqref{eq:geometry}. At all later
coordinates, the sum over \emph{every} internal layer, including the
pivot, is at most
\[
      \sigma(1+1/7)\sum_{q\ge1}B^{-q}=\frac{8}{105}\sigma .
\]
Both bounds hold on every path for every setting of the unexposed bits.
With $\rho=1/7+8/105=23/105<1/4$, this gives
\begin{equation}\label{eq:sandwich}
  c_u+\sigma\eta_j a_j(u)\ \le\ D_u\
            \le\ c_u+\sigma\eta_j a_j(u)+\rho\sigma .
\end{equation}
Nonnegativity gives the lower bound. For the upper bound, evaluate the
full cost on any path attaining $c_u$. This argument allows the
minimizing path to change with the unexposed bits.

Applying \eqref{eq:sandwich} to both classes shows that $D_u=D_v$
requires
\begin{equation}\label{eq:forced}
   \left|\eta_j\bigl(a_j(u)-a_j(v)\bigr)
              -\frac{c_v-c_u}{\sigma}\right|\le\rho .
\end{equation}
When the word bits differ, the two values on the left before subtracting
the fixed center are integers a unit apart. At most one choice of
$\eta_j$ can satisfy \eqref{eq:forced}, because $2\rho<1$.

Fix all mask bits at coordinates where the two words agree. If two
completions both gave a tie, take their first differing mask bit.
The prefixes through $j-1$ agree, so the center in \eqref{eq:forced}
is the same for both completions, while the two choices of $\eta_j$
cannot both satisfy it. Thus at most one completion ties the two
class minima. There are $2^{\dist(a(u),a(v))}$ equiprobable
completions, proving the bound.
\end{proof}

\section{Simultaneous isolation and unambiguous search}

\begin{proof}[Proof of Theorem~\ref{thm:main}]
For every fixed distinct pair $u,v$, the difference
$a(u)\oplus a(v)$ is uniform on $\F_2^k$. Indeed, let $i_*$ be the
largest position where $x(u)$ and $x(v)$ differ. In row $j$ the
difference uses $\tau_{j-1+i_*}$ with coefficient one, and its other
effective indices are smaller. This bit is fresh relative to rows
$1,\ldots,j-1$. The difference bits are therefore independent and uniform.

Fix an endpoint pair and two feasible pivot vertices $u,v$. By
Lemma~\ref{lem:mask}, averaging over the word table gives
\[
   \Pr_{\tau,\eta}[D_u=D_v]\le
    \mathbb E_\tau[2^{-\dist(a(u),a(v))}]
      =\left(\tfrac12+\tfrac12\cdot\tfrac12\right)^k=(3/4)^k.
\]
There are at most $N^4$ choices of the two endpoints and two pivot
vertices. The probability of any such equality is at most
$N^4(3/4)^{16L}\le N^{-2}$, since $(3/4)^{16}<2^{-6}$.

Fix a seed outside this event. We prove uniqueness for every reachable
pair by induction on layer separation. Pairs with no internal layer have
a unique path by simplicity, including the trivial path from a vertex
to itself. Otherwise, among feasible pivot vertices the values $D_v$
are all distinct. Every minimum path uses the unique winning pivot.
The subpaths on its two sides must each be minimum between their own
endpoints, by replacement. Both are unique by induction.

The seed has $2k+L-1$ bits. Every vertex weight is less than
$B^{L+k}$, and any path uses at most $N$ vertices. To compute a weight
one evaluates the word bits by XOR and uses binary shifts, since
$B=16$. This takes $O(L)$ space.
\end{proof}

\begin{corollary}\label{cor:ul}
$\NL=\UL$.
\end{corollary}
\begin{proof}
Layered directed reachability is $\NL$-complete under deterministic
logspace reductions~\cite[Proposition 3.4]{vMP}. One may time-unroll a
logspace configuration graph for polynomially many steps, allow padding
after acceptance, and merge parallel edges. The resulting graph has
polynomially many vertices, and adjacency can be queried in logspace.

Enumerate the seeds of Theorem~\ref{thm:main} in a fixed order.
Use a total unambiguous check for min-isolation on each seed, and
decide reachability with the first passing assignment. This is the
isolation-to-unambiguity reduction of Reinhardt and Allender~\cite{RA}
in the weighted form of van Melkebeek and Prakriya~\cite[Lemma 3.5
and Section 3.2]{vMP}. Their space bound is $O(\log N+\log(W+1))$
for an upper bound $W$ on path weights. One can add one to every
vertex weight if positive weights are required; the number of vertices
on a path is fixed by its endpoints. The seed counter, weight evaluation,
and the check use $O(\log N)$ space, hence $O(\log n)$ for original
input length $n$. There are polynomially many seeds. Theorem~\ref{thm:main}
ensures a passing seed, and $\UL\subseteq\NL$ follows from the definition.
\end{proof}

\paragraph{Why failed seeds preserve unambiguity.}
Here is a direct implementation of the check on layered graphs. Charge
the vertices entered by a path, including its last vertex. For layer
separation $j$, keep the number of reachable ordered pairs and the sum
of their minimum distances. For $j=0$ the census is $(N,0)$.
Given a correct census with unique minima at separation $j$, a lookup
of any one distance has exactly one returning computation. Scan all
pairs at that separation in fixed order. For each pair either skip it
or guess a path between its endpoints, using a unique encoding of the
vertex sequence. Abort unless the number of paths and their total
weight match the census. The count forces a path for every reachable
pair; the sum forces each path to be minimum. This also yields the
requested distance, with infinity for an unreachable pair.

For each pair at separation $j+1$, scan its possible last predecessors
and use these lookups. Two different predecessors attaining a finite
minimum certify failure of the seed, with a unique returning computation.
If this never happens, accumulate the census for $j+1$ and continue.
The designated reachability answer can be recorded on the way. The scan
uses a constant number of $O(\log N)$-bit counters and vertex names.
Thus a failing seed has a unique computation continuing to the next seed,
and a passing seed has a unique computation returning the answer.

\paragraph{Dependence on the graph.}
The word and mask family itself is oblivious to the edges. The successful
seed selected by the check depends on the particular graph. Hoza's
lower bound concerns a single assignment that must work for all graphs
of given width and length simultaneously~\cite{Hoza}.

\begin{thebibliography}{9}
\bibitem{Hoza} William M. Hoza. \emph{Deterministic, Oblivious Isolation for
Space-Bounded Computation Requires Large Weights}. ECCC TR26-124, 2026.
\url{https://eccc.weizmann.ac.il/report/2026/124/}
\bibitem{RA} Klaus Reinhardt and Eric Allender. \emph{Making Nondeterminism
Unambiguous}. SIAM Journal on Computing 29(4), 1118--1131, 2000.
\href{https://doi.org/10.1137/S0097539798339041}{doi:10.1137/S0097539798339041}.
\bibitem{vMP} Dieter van Melkebeek and Gautam Prakriya. \emph{Derandomizing
Isolation in Space-Bounded Settings}. SIAM Journal on Computing 48(3),
979--1021, 2019.
\href{https://doi.org/10.1137/17M1130538}{doi:10.1137/17M1130538}.
\end{thebibliography}
\end{document}
